\documentclass[10pt,a4paper]{article}
\usepackage[T1]{fontenc}
\usepackage[utf8]{inputenc}
\usepackage{mathptmx}
\usepackage{courier}
\usepackage[english]{babel}
\usepackage{amsmath,amssymb}
\usepackage[margin=1.9cm]{geometry}
\usepackage{microtype}
\usepackage{booktabs,tabularx,array}
\usepackage{titlesec}
\usepackage{caption}
\pagestyle{empty}
\titleformat{\section}{\normalsize\bfseries\scshape}{\thesection.}{0.5em}{}
\titlespacing*{\section}{0pt}{1.1ex plus .3ex}{0.5ex}
\setlength{\parskip}{0.3ex}
\renewcommand{\arraystretch}{1.08}
\newcolumntype{Y}{>{\raggedright\arraybackslash}X}
\captionsetup{font=small,labelfont=bf,skip=3pt}
\setlength{\textfloatsep}{8pt}
\setlength{\intextsep}{6pt}

\begin{document}
\begin{center}
{\large\bfseries Does Underlying Structure Predict Quantum Trainability and Dequantization?}\\[0.3ex]
{\small DALAS, Deliverable 1 (TME4): research question and dataset}\\[0.2ex]
{\small Group: Charbel A. Dargham (solo, scope approved by the instructor) \quad\textperiodcentered\quad TME1}
\end{center}

\section{Research question}

\emph{Can the underlying structure of a learning problem predict, before training, whether structure-preserving
quantum models of it will be trainable and resist dequantization?} Structure-free quantum circuits become
untrainable as they grow, their gradients vanishing exponentially with the number of qubits (McClean et al.\
2018); building a problem's symmetry into the circuit keeps them trainable (Larocca et al.\ 2022), but can
also let a classical computer imitate the model (Cerezo et al.\ 2025). Which side a problem lands on is today
known only after a model is built and run. We ask whether it can be read off the problem itself, from four
kinds of structure: symmetry, locality, spectrum and data complexity. This deliverable builds the first
ingredient, a large and varied collection of real learning problems with their structure extracted.

\section{Data collection}

The unit of observation is a \emph{problem}: one binary learning task on $n$ sites, with its labelled inputs,
a TRAIN/TEST split and its label rule. Each input $x$ is stored with an operator $M(x)$, and the operator fixes
the problem's type (Table~\ref{tab:types}): a graph Laplacian defines a heat equation on the graph, a precision
matrix (inverse covariance) an Ornstein--Uhlenbeck process, and a Hamiltonian a Schr\"odinger equation. The type
depends on the mathematics, not the domain: a MaxCut instance encodes like a graph, while a 3-SAT formula has
three-site terms a graph cannot have. Sizes run from $n = 4$, below which almost every structure is trivial, to
$n = 12$, where states still have only $2^{12} = 4096$ amplitudes and exact simulation stays cheap.

\begin{table}[h]
\centering\small
\caption{Types, sources and counts. Groups are independent sources (one dataset, network or generator).}\label{tab:types}
\begin{tabularx}{\linewidth}{@{}l l Y r r r@{}}
\toprule
Type & Operator $M(x)$ & Sources and collection & Problems & Groups & Samples\\
\midrule
GRAPH & graph Laplacian & 20 real networks (Netzschleuder API: biology, social, infrastructure, technology), cut into connected pieces of $n$ nodes; real molecules (Alchemy, ZINC via TUDataset) & 345 & 22 & 100\,782\\
COVARIANCE & precision matrix & EEG motor imagery (MOABB API: BNCI2014\_001, Cho2017), one problem per subject and session; 9 sensor datasets (UEA archive) & 404 & 11 & 87\,992\\
DIAGONAL & clause-penalty Hamiltonian & SAT formulas from 11 standard generators (random $k$-SAT, XOR, NAE, planted, community, SR) & 358 & 11 & 107\,400\\
QUANTUM & Hamiltonian, ground state & 8 spin-chain models (exact ground states on both sides of a phase transition); 6 molecules (Hartree--Fock, Jordan--Wigner) & 261 & 14 & 41\,756\\
\midrule
Total & & & 1\,368 & 58 & 337\,930\\
\bottomrule
\end{tabularx}
\end{table}

Each source was listed, screened, downloaded (with a file hash) and converted by a logged pipeline; every
rejection is recorded with its reason (for the real networks, 892 candidates were checked and 307 kept). A
candidate problem was kept only if it had at least 100 samples, so that TEST holds enough samples to measure an
accuracy; at least 50 distinct inputs, since a benchmark such as IMDB-MULTI at $n = 7$ holds 332 graphs but only
10 different ones; and at most 20\% of equal inputs carrying different labels, beyond which the label is no
longer a function of the input. Labels are either given by the source or computed exactly from the input; when a
source gives a score, the label is $+1$ above the TRAIN median, so the TEST data never influence the cut. In 51
problems (42 COVARIANCE, 9 GRAPH) the share of $+1$ labels fell outside $[0.2, 0.8]$, and the majority class was
randomly reduced to the size of the minority class; each problem records what was done. The split is 70/30 by
convention; it is stratified (same class proportions on both sides) and made by unit, so that samples derived
from one original (two crops of one EEG trial) never straddle it. Real inputs make up 749 problems (55\%); the
SAT and spin-chain problems are generated from standard models, because in these domains generated instances
are how problems are produced, and every problem carries its origin.

\section{Variables}

The variables are a profile of about a dozen numbers per input, grouped under the four structures. A problem is
then a cloud of such profiles, one per input, and how to summarise that cloud is an open question, deferred to
the analysis. Each problem records its type, $n$, number of samples, split, class balance, label rule, number of distinct
inputs, conflict rate, source, group, origin, and an equation check (0 failures in all 1\,368 problems). Each
sample records an input $x$, its operator $M(x)$ and its label. Type, label rule, source, group, origin and
the label ($\pm1$) are categorical; $n$, the counts, the rates and every structure variable are numerical. The structure variables (Table~\ref{tab:vars})
are computed \emph{per input}, from the \emph{factor graph} of $M(x)$: a bipartite graph with one node per site
and one per term, a term joined to every site it acts on. It is the standard site-level object for a sum of
terms (Kschischang et al.\ 2001), and unlike a pairwise coupling matrix it keeps many-body terms apart from
pairwise ones. Weights are divided by the largest, so no variable depends on units.

\begin{table}[h]
\centering\small
\caption{Structure variables, one value per input, grouped by the four kinds of structure.}\label{tab:vars}
\begin{tabularx}{\linewidth}{@{}l >{\raggedright\arraybackslash}p{5.2cm} Y@{}}
\toprule
Structure & Variable & Meaning\\
\midrule
Symmetry & symmetry classes, fraction & how many sites look alike under colour refinement (1: all alike, 0: none); exact equality of weights, and an upper bound on true symmetry (Grohe et al.\ 2014)\\
Locality & terms per site; term sizes & interactions per site; whether terms are pairwise or many-body\\
 & site degree spread; support density & whether a few sites dominate; how far interactions reach\\
Spectrum & spectral entropy; participation ratio & how many Laplacian modes of the factor graph carry its weight\\
 & algebraic connectivity & how hard the factor graph is to cut in two\\
 & operator spectrum (GRAPH, COVARIANCE) & entropy and gap ratio of the eigenvalues of $M(x)$ itself\\
Data complexity & entanglement (QUANTUM); effective rank (COVARIANCE) & half-chain entropy of the ground state; number of effectively independent channels\\
\bottomrule
\end{tabularx}
\end{table}

The factor-graph spectrum is a proxy: it describes which sites the terms connect, while the spectrum of $M(x)$
itself is available for GRAPH and COVARIANCE only. Data complexity is defined on the encoded quantum state of
each input; for GRAPH and DIAGONAL it is computed at the encoding stage, after this deliverable. The
type-specific variables are therefore missing by design wherever they do not exist (entanglement outside
QUANTUM, effective rank outside COVARIANCE, the operator spectrum outside GRAPH and COVARIANCE): the missingness
is set by the type, not by chance, so these variables are compared within a type only.

\emph{The dataset as a whole.} The types occupy different regions of this space, partly by construction. GRAPH
operators have pairwise terms only, and their density and symmetry vary with the real network they come from.
COVARIANCE operators connect every pair of channels, so their locality is uniform and their symmetry is zero by
exact equality: their structure lies in the weights, the partial correlations. DIAGONAL operators carry terms of
up to $k$ sites, set by the generator, so their structure is predictable from its parameters. Spin chains have
one- and two-site terms (three for the cluster model) and full symmetry when periodic; molecules have terms
spanning up to all $n$ sites.
\section{Preprocessing and choices}

The raw harvest (1\,379 problems) is never modified; a script with six stated rules builds the final dataset
from it, and the raw files' checksum is identical before and after. An audit found 0 structural failures and
five problems in the content, each handled by one choice. Identical inputs on both sides of the split, in 22
problems, would leak TEST into TRAIN: these problems were re-split with identical inputs kept on one side,
rather than deduplicated, which would have changed the population, and repeated inputs are weighted so that
each counts once in TEST. One problem below the class-balance floor was excluded. In 24\% of spin-chain
samples the ground state was degenerate, so the stored state was arbitrary: the spin chains were re-harvested
with ground states fixed inside a symmetry sector. Two molecule geometries whose Hartree--Fock calculation did
not converge were removed. Finally, 475 sensor samples whose operator collapses to the identity carry no
information and act as a label shortcut; they are flagged rather than removed, and every result is reported with
and without them. The periodic XY spin chains (10 candidate problems) are not part of the dataset: the
eigensolver failed on one sample whose excited levels nearly coincide, and the pipeline rejected the whole
family, as designed. The final dataset passes the audit with 0 failures and 0 leaks.

\section{Conclusion and next steps}

The dataset gathers 1\,368 binary learning problems of four mathematical types, each input stored with the
operator that defines its equation and each step from source to final file logged and checked. Its strengths
are its breadth (58 independent sources, 55\% of problems on real inputs) and its traceability: every rejection,
re-split, exclusion and flag can be read back. Its limits are stated with it. The effective sample size is set by
the groups, not the problems, and it is uneven: 350 of the 404 COVARIANCE problems come from two EEG datasets.
The DIAGONAL and most QUANTUM problems are generated from standard models. Data complexity is measured for two
types only so far. The next steps follow the chain of the research question. First, decide how the per-input
values of one problem are summarised into one profile, after plotting their distributions to check whether they
form one group or several within a problem, which a single summary would hide. Second, build a first instance of a Logical State Space $L = (C; F; \Phi)$, the
construction of the author's ongoing research, intended as the bridge between a problem and its circuit. A
predictor alone can only show that structure and outcome are correlated; $L$ gives every problem one common
medium in which that correlation can be explained in mathematical terms, as connectomics relates a brain's
behaviour to its wiring, a link that for model neural networks is a theorem (Curto, Geneson \& Morrison 2019). In this first instance, $C$ is
the simplicial complex generated by the terms of each factor graph, $F$ the constant sheaf over it, and $\Phi$
diffusion under its Hodge Laplacians, whose non-decaying modes count the holes of $C$ (Eckmann 1944), a theorem
that ties structure to dynamics. Circuits are then generated from $L$, each of the four structures entering as a
separate ingredient, so that circuits built with and without each one measure its effect; the inputs are
encoded and the data-complexity variables completed. Third, simulate the circuits and measure trainability, through gradient variance and the dimension of
the dynamical Lie algebra (Ragone et al.\ 2024), and classical simulability (Cerezo et al.\ 2025). Last, learn
the predictor from structure to these measurements, validated by holding out whole sources.

\end{document}
